\documentclass[11pt]{article}
\usepackage[a4paper,margin=20mm]{geometry}
\usepackage{fontspec}
\setmainfont{Noto Serif CJK SC}
\XeTeXlinebreaklocale "zh"
\XeTeXlinebreakskip=0pt plus 1pt
\usepackage{amsmath,amssymb,amsthm,mathtools,tikz-cd,enumerate,paralist,xurl,hyperref}
\setlength{\emergencystretch}{4em}
\newcommand{\identity}{\mathrm{id}}
\newcommand{\Obj}{\operatorname{Ob}}
\newcommand{\Mor}{\operatorname{Mor}}
\newcommand{\Hom}{\operatorname{Hom}}
\newcommand{\End}{\operatorname{End}}
\newcommand{\Ker}{\operatorname{ker}}
\newcommand{\Coker}{\operatorname{coker}}
\newcommand{\opp}{\mathrm{op}}
\newcommand{\munit}{\mathbf{1}}
\newcommand{\Z}{\mathbb{Z}}
\newcommand{\cate}[1]{\mathbf{#1}}
\newcommand{\cated}[1]{\mathbf{#1}\text{-}}
\newcommand{\dcate}[1]{\text{-}\mathbf{#1}}
\newcommand{\dotimes}[1]{\mathbin{\mathop{\otimes}\limits_{#1}}}
\newcommand{\rightiso}{\xrightarrow{\sim}}
\renewcommand{\index}[2][]{\relax}
\newcommand{\CorpusNumber}{0}
\newtheorem{theorem}{定理}
\newtheorem{lemma}[theorem]{引理}
\newtheorem{proposition}[theorem]{命题}
\newtheorem{corollary}[theorem]{推论}
\theoremstyle{definition}
\newtheorem{definition}[theorem]{定义}
\newtheorem{definition-proposition}[theorem]{定义--命题}
\newtheorem{example}[theorem]{例}
\newtheorem{remark}[theorem]{注记}
\renewcommand{\thetheorem}{\CorpusNumber}
\renewcommand{\proofname}{证明}
\begin{document}
\section*{294: Morita equivalence classification}
\noindent\textbf{作者：} 李文威 (Wen-Wei Li). \textbf{作品：}《代数学方法》第二卷，第七章。
\par\noindent\textbf{来源与版本：} \url{https://github.com/wenweili/AlJabr-2/blob/e03224e2e867a49dc60c7975d890079068a30466/chapter7.tex#L1977}, commit \texttt{e03224e2e867a49dc60c7975d890079068a30466}.
\par\noindent\textbf{许可：} Creative Commons Attribution 4.0 International; copyright 2024 李文威. \url{https://creativecommons.org/licenses/by/4.0/}. 作者不保证无误；见附带 LICENSE。
\par\noindent\textbf{整理说明（非作者正文）：} 下列题面、证明及局部支撑正文来自作者原生 TeX；保留原语言、顺序及数学内容。仅改独立排版、原编号引用与索引命令；未生成或补写证明。所有文字与数学图表均可编辑。标准先修仅引用，不递归重证。
\par\noindent\textbf{难度说明（整理者）：} H2: The proof reconstructs bimodules implementing an equivalence and proves evaluation and coevaluation are isomorphisms from the finite projective generator hypotheses. The core Morita equivalence is proved rather than assumed; Eilenberg-Watts is recorded prerequisite material.
\subsection*{原命题与完整人类证明}
\def\CorpusNumber{7.7.16}
\begin{theorem}[森田纪一]\label{prop:Morita-refined}
	考虑 $\Bbbk$-代数 $A$ 和 $B$. 命 $\mathrm{Equiv}(\cated{Mod}A, \cated{Mod}B)$ 为以所有等价 $F: \cated{Mod}A \to \cated{Mod}B$ 为对象, 以其间的同构为态射的范畴. 另一方面, 定义范畴 $\mathcal{P}(A, B)$ 使得其对象是满足下述条件的 $(A, B)$-双模 $P$:
	\begin{itemize}
		\item $P$ 作为右 $B$-模是 $\cated{Mod}B$ 的有限生成投射生成元,
		\item 左乘诱导同构 $A \rightiso \End_{\cated{Mod}B}(P)$,
	\end{itemize}
	其态射则定为双模的同构. 我们有以下互为拟逆的函子:
	\[\begin{tikzcd}[row sep=tiny]
		\mathrm{Equiv}(\cated{Mod}A, \cated{Mod}B) \arrow[shift left, r] & \mathcal{P}(A, B) \arrow[shift left, l] \\
		F \arrow[mapsto, r] & F(A) \\
		(\cdot) \dotimes{A} P & P \arrow[mapsto, l] .
	\end{tikzcd}\]
\end{theorem}
\begin{proof}
	首先说明 $F \mapsto F(A)$ 是良定义的. 观察到 $A$ 是 $\cated{Mod}A$ 的有限生成投射生成元, 既然投射生成元和有限生成性质皆有范畴论刻画 (引理 7.7.15), 故 $F(A)$ 亦然; 此外, 根据抽象的理由,
	\[ F: \End_{\cated{Mod}A}(A) \rightiso \End_{\cated{Mod}B}(F(A)), \]
	不难看出这正是 $A$ 对 $F(A)$ 的左乘给出的 $A \to \End_{\cated{Mod}B}(F(A))$.
	
	若能说明从右到左的函子也是良定义的, 则定理 7.7.2 便蕴涵双向的函子互为拟逆. 问题遂归结为证: 若 $P \in \Obj(\mathcal{P}(A, B))$, 则 $F := (\cdot) \dotimes{A} P$ 是范畴的等价.
	
	取 $P$ 如上. 在 (7.7.6) 已经说明 $(\cdot) \dotimes{A} P$ 的右伴随是 $(\cdot) \dotimes{B} P^\vee$, 其中 $P^\vee := \Hom_{\cated{Mod}B}(P, B)$. 目标是说明它们互为拟逆. 在定义--命题 7.7.4 (ii) 中取 $Q = P^\vee$, 将问题进一步化约为证明定义 7.7.10 的双模同态
	\[ \mathrm{ev}: P^\vee \dotimes{A} P \to B, \quad \mathrm{coev}: A \to P \dotimes{B} P^\vee \simeq \End_{\cated{Mod}B}(P) \]
	皆为同构.
	
	先看 $\mathrm{ev}$. 由 $P$ 是 $\cated{Mod}B$ 的生成元可知存在满同态 $P^{\oplus I} \twoheadrightarrow B$, 特别地, 存在 $q_1, \ldots, q_n \in P^\vee$ 和 $p_1, \ldots, p_n \in B$ 使得 $\sum_{i=1}^n q_i(p_i) = 1_B$, 亦即 $\mathrm{ev}(\sum_i q_i \otimes p_i) = 1_B$. 满性得证.
	
	接着引进写作乘法的运算 $P \times P^\vee \to A$, 它的刻画是对于所有 $p, p' \in P$ 和 $q \in P^\vee$, 在 $P$ 中有形似结合律的
	\[ \underbracket{(pq)}_{\in A} p' = p \underbracket{(qp')}_{\in B}, \]
	其中 $qp' := q(p')$; 诚然, 当 $p$ 和 $q$ 固定, 右式是 $\End_{\cated{Mod}B}(P)$ 的元素, 而左乘诱导 $A \rightiso \End_{\cated{Mod}B}(P)$, 由此唯一定义了 $pq \in A$. 从上式还可以推导第二种结合律, 这是 $P^\vee$ 中的等式
	\[ \underbracket{(qp)}_{\in B} q' = q \underbracket{(pq')}_{\in A}; \]
	说明两边对所有 $p' \in P$ 取值皆相等即可, 平凡的验证留给读者.

	现在可以说明 $\mathrm{ev}$ 的单性. 选定上一步得到的 $p_i$ 和 $q_i$, 设 $\mathrm{ev}(\sum_{j=1}^m q'_j \otimes p'_j) = 0$, 则上述两种结合律蕴涵
	\begin{multline*}
		\sum_j q'_j \otimes p'_j = \sum_{i, j} (q_i p_i) q'_j \otimes p'_j
		= \sum_{i, j} q_i \underbracket{(p_i q'_j)}_{\in A} \otimes p'_j \\
		= \sum_{i, j} q_i \otimes (p_i q'_j) p'_j
		= \sum_{i, j} q_i \otimes p_i \underbracket{(q'_j p'_j)}_{\in B}
		= \sum_j q_i \otimes p_i \sum_j q'_j p'_j = 0.
	\end{multline*}

	最后, 关于 $\mathrm{coev}$ 为同构的断言无非是 $\mathcal{P}(A, B)$ 定义的第二部分.
\end{proof}
\subsection*{必要作者背景与局部支撑（不另计题）}
\par\medskip\noindent\textit{作者原文：chapter7.tex，行 1662--1688.}\par
本节选定 Grothendieck 宇宙 $\mathcal{U}$ 和交换环 $\Bbbk$. 相对于 $\mathcal{U}$, 所有的环都默认是小的. 对于 $\Bbbk$-线性范畴之间的函子同样默认 $\Bbbk$-线性, 如定义 1.4.1, 1.4.4 所述. 当 $\Bbbk = \Z$ 时, $\Bbbk$-线性也就是加性.

对于 $\Bbbk$-线性范畴 $\mathcal{C}$ 和 $\mathcal{D}$, 符号 $\mathrm{Fct}(\mathcal{C}, \mathcal{D})$ 或 $\mathcal{D}^{\mathcal{C}}$ 代表所有函子 $F: \mathcal{C} \to \mathcal{D}$ 构成的范畴. 尽管 $\mathrm{Fct}(\mathcal{C}, \mathcal{D})$ 未必是 $\mathcal{U}$-范畴, 由于我们不在其中作极限等构造, 这不会对往后的讨论造成任何困难.

\def\CorpusNumber{7.7.1}
\begin{definition}
	\index[sym1]{Fctc@$\mathrm{Fct}^c, \mathrm{Fct}_c$}
	设 $\mathcal{C}$ 和 $\mathcal{D}$ 为 $\Bbbk$-线性范畴.
	\begin{itemize}
		\item 设 $\mathcal{C}$ 和 $\mathcal{D}$ 余完备, 定义 $\mathrm{Fct}(\mathcal{C}, \mathcal{D})$ 的全子范畴 $\mathrm{Fct}^c(\mathcal{C}, \mathcal{D})$, 由全体保小 $\varinjlim$ 的函子组成;
		\item 设 $\mathcal{C}$ 和 $\mathcal{D}$ 完备, 定义 $\mathrm{Fct}(\mathcal{C}, \mathcal{D})$ 的全子范畴 $\cate{Fct}_c(\mathcal{C}, \mathcal{D})$, 由全体保小 $\varprojlim$ 的函子组成.
	\end{itemize}
	一些文献称 $\mathrm{Fct}_c$ (或 $\mathrm{Fct}^c$) 的对象为连续 (或余连续) 函子.
\end{definition}

我们沿用 \S3.10 的符号, 对任意 $\Bbbk$-代数 $A$ 和 $B$, 定义
\begin{align*}
	(A, B)\dcate{Mod} & := (A, B)\text{-双模范畴}, \\
	A\dcate{Mod} & := \text{左}\; A\text{-模范畴} \simeq (A, \Bbbk)\dcate{Mod}, \\
	\cated{Mod}B & := \text{右}\; B\text{-模范畴} \simeq (\Bbbk, B)\dcate{Mod}.
\end{align*}
它们是 $\Bbbk$-线性范畴的简单例子.

回忆以下的模论常识: $(A, B)$-双模 $P$ 典范地诱导保小 $\varinjlim$ 的函子
\[ P \dotimes{B} (\cdot): B\dcate{Mod} \to A\dcate{Mod}, \quad (\cdot) \dotimes{A} P : \cated{Mod}A \to \cated{Mod}B, \]
而 $(B, A)$-双模 $Q$ 典范地诱导保小 $\varprojlim$ 的函子
\[ \Hom_{\cated{Mod}B}(Q, \cdot): B\dcate{Mod} \to A\dcate{Mod}, \quad \Hom_{\cated{Mod}A}(Q, \cdot): \cated{Mod}A \to \cated{Mod}B, \]
若改取函子 $\Hom(\cdot, R)$, 其中 $R$ 是 $(B, A)$-双模, 则其定义域须换为相反范畴如 $(B\dcate{Mod})^{\opp}$ 等, 使函子依然保小 $\varprojlim$.
\par\medskip\noindent\textit{作者原文：chapter7.tex，行 1690--1702.}\par
\begingroup\small
\def\CorpusNumber{7.7.2}
\begin{theorem}[S.\ Eilenberg, C.\ E.\ Watts]\label{prop:Morita}
	对于任意 $\Bbbk$-代数 $A$ 和 $B$, 我们有以下等价
	\[\begin{tikzcd}[row sep=tiny]
		\mathrm{Fct}^c\left(B\dcate{Mod}, A\dcate{Mod}\right) & (A, B)\dcate{Mod} \arrow[l, "\sim"'] \arrow[r, "\sim"] & \mathrm{Fct}^c\left(\cated{Mod}A, \cated{Mod}B\right) \\
		P \dotimes{B} (\cdot) & P \arrow[mapsto, l] \arrow[mapsto, r] & (\cdot) \dotimes{A} P \\
		\mathrm{Fct}_c\left(B\dcate{Mod}, A\dcate{Mod}\right) & (B, A)\dcate{Mod} \arrow[l, "\sim"'] \arrow[r, "\sim"] & \mathrm{Fct}_c\left(\cated{Mod}A, \cated{Mod}B\right) \\
		\Hom_{B\dcate{Mod}}(Q, \cdot) & Q \arrow[mapsto, l] \arrow[mapsto, r] & \Hom_{\cated{Mod}A}(Q, \cdot) \\
		\mathrm{Fct}_c\left((B\dcate{Mod})^{\opp}, \cated{Mod}A \right) & (B, A)\dcate{Mod} \arrow[l, "\sim"'] \arrow[r, "\sim"] & \mathrm{Fct}_c\left((\cated{Mod}A)^{\opp}, B\dcate{Mod}\right) \\
		\Hom_{B\dcate{Mod}}(\cdot, R) & R \arrow[mapsto, l] \arrow[mapsto, r] & \Hom_{\cated{Mod}A}(\cdot, R).
	\end{tikzcd}\]

	对于第一和第二个等价, 函子的合成对应到双模的张量积, 而 $A=B$ 时恒等函子来自作为 $(A, A)$-双模的 $A$, 精确到同构.
\end{theorem}
\endgroup
\par\medskip\noindent\textit{作者原文：chapter7.tex，行 1736--1752.}\par
\def\CorpusNumber{7.7.4}
\begin{definition-proposition}\label{def:Morita-equiv}
	\index{sentiandengjia@森田等价 (Morita equivalence)}
	对于任意 $\Bbbk$-代数 $A$ 和 $B$, 以下陈述等价.
	\begin{enumerate}[(i)]
		\item $A\dcate{Mod}$ 和 $B\dcate{Mod}$ 等价.
		\item 存在 $(A, B)$-双模 $P$ 和 $(B, A)$-双模 $Q$ 使得存在
		\begin{compactitem}
			\item $(A, A)$-双模的同构 $P \dotimes{B} Q \simeq A$,
			\item $(B, B)$-双模的同构 $Q \dotimes{A} P \simeq B$.
		\end{compactitem}
		\item $\cated{Mod}A$ 和 $\cated{Mod}B$ 等价.
	\end{enumerate}
	当以上任一条件成立时, 我们称 $A$ 和 $B$ 是\emph{森田等价}的.
\end{definition-proposition}
\begin{proof}
	等价必然保 $\varinjlim$ 和 $\varprojlim$. 从定理 7.7.2 可得 (i) $\iff$ (ii) 和 (iii) $\iff$ (ii).
\end{proof}
\par\medskip\noindent\textit{作者原文：chapter7.tex，行 1773--1786.}\par
\begingroup\small
回到一般情形. 对任意 $(A, B)$-双模 $P$, 模论的常识 \cite[定理 6.6.5]{Li1} 给出伴随对
\renewcommand{\theequation}{7.7.2}
\begin{equation}\label{eqn:AB-adjunction}\begin{tikzcd}
	(\cdot) \dotimes{A} P : \cated{Mod}A \arrow[shift left, r] & \cated{Mod}B: \Hom_{\cated{Mod}B}(P, \cdot). \arrow[shift left, l]
\end{tikzcd}\end{equation}
定理 7.7.2 说明精确到同构, 左右两端的函子分别穷尽了 $\mathrm{Fct}^c(\cated{Mod}A, \cated{Mod}B)$ 和 $\mathrm{Fct}_c(\cated{Mod}B, \cated{Mod}A)$ 的所有对象.

基于例 7.6.2, 伴随对 (7.7.2) 在 $\cated{Mod}A$ 上确定单子 $T$, 在 $\cated{Mod}B$ 上确定余单子 $L$, 其一般描述如下. 考虑右 $A$-模 $N$ 和右 $B$-模 $M$:
\renewcommand{\theequation}{7.7.3}
\begin{equation}\label{eqn:bimod-adjunction}
	\begin{array}{|c|c|} \hline
		\text{单位}\; \eta_N & \text{余单位}\; \epsilon_M \\ \hline
		N \to \Hom_{\cated{Mod}B}\left( P, N \dotimes{A} P \right) & \Hom_{\cated{Mod}B}(P, M) \dotimes{A} P \to M \\
		x \mapsto \left[ p \mapsto x \otimes p \right] & \varphi \otimes p \mapsto \varphi(p) \\ \hline
	\end{array}
\end{equation}
\endgroup
\par\medskip\noindent\textit{作者原文：chapter7.tex，行 1831--1888.}\par
其次考虑一个稍广的情形. 首先, 对任何右 $B$-模 $P$ 和 $X$, 视 $B$ 为 $(B, B)$-双模以定义左 $B$-模
\renewcommand{\theequation}{7.7.4}
\begin{equation}\label{eqn:P-vee}
	P^\vee := \Hom_{\cated{Mod}B}(P, B)
\end{equation}
以及典范的 $\Bbbk$-模同态
\renewcommand{\theequation}{7.7.5}
\begin{equation}\label{eqn:dual-Hom-X}\begin{aligned}
	X \dotimes{B} P^\vee & \to \Hom_{\cated{Mod}B}(P, X) \\
	x \otimes \lambda & \mapsto x \lambda(\cdot).
\end{aligned}\end{equation}
\index[sym1]{Pvee@$P^\vee$}

如果 $P$ 是 $(A, B)$-双模, 则 $P^\vee$ 具有 $(B, A)$-双模的结构, $\Hom_{\cated{Mod}B}(P, X)$ 具有右 $A$-模结构, 而 (7.7.5) 是右 $A$-模同态.

相同构造当然也有左右对调的版本, 故 $P^{\vee\vee}$ 有意义, 此外还有右 $B$-模的典范同态 $P \to P^{\vee\vee}$.

\def\CorpusNumber{7.7.9}
\begin{lemma}\label{prop:AB-comonad-aux}
	设 $(A, B)$-双模 $P$ 作为右 $B$-模是有限生成投射模, 则 $P^\vee$ 亦然, 此时
	\begin{itemize}
		\item 对所有 $X$, (7.7.5) 皆给出右 $A$-模的典范同构,
		\item $P \to P^{\vee\vee}$ 是同构.
	\end{itemize}
\end{lemma}
\begin{proof}
	将 $P$ 作为右 $B$-模表作 $B^{\oplus n}$ 的直和项, $n \in \Z_{\geq 0}$. 于是 $P^\vee$ 自然地是 $B^{\oplus n}$ 的直和项. 显见在 (7.7.5) 中以 $B^{\oplus n}$ 代 $P$ 可得 $\Bbbk$-模同构, 于是在直和项 $P$ 的层次仍有 $\Bbbk$-模同构, 因而有右 $A$-模同构. 关于 $P \rightiso P^{\vee\vee}$ 同样是化到 $B^{\oplus n}$ 来论证.
\end{proof}

于是伴随对 (7.7.2) 在上述假设下改写为
\renewcommand{\theequation}{7.7.6}
\begin{equation}\label{eqn:AB-adjunction-2}\begin{tikzcd}
	(\cdot) \dotimes{A} P : \cated{Mod}A \arrow[shift left, r] & \cated{Mod}B: (\cdot) \dotimes{B} P^\vee . \arrow[shift left, l]
\end{tikzcd}\end{equation}

\def\CorpusNumber{7.7.10}
\begin{definition}\label{def:bimodule-ev-coev}
	\index[sym1]{ev-coev@$\mathrm{ev}, \mathrm{coev}$}
	设 $(A, B)$-双模 $P$ 作为右 $B$-模是有限生成投射模. 定义 $(B, B)$-双模同态
	\[\begin{tikzcd}[row sep=tiny]
		\mathrm{ev}: & P^\vee \dotimes{A} P \arrow[r] & B \\
		& \lambda \otimes p \arrow[mapsto, r] & \lambda(p)
	\end{tikzcd}\]
	和 $(A, A)$-双模同态
	\[\begin{tikzcd}[row sep=tiny]
		\mathrm{coev}: & A \arrow[r] & P \dotimes{B} P^\vee \arrow[r, "{\text{引理 7.7.9}}" inner sep=0.6em, "\sim"'] & \End_{\cated{Mod}B}(P) \\
		& a \arrow[mapsto, rr] & & {[p \mapsto ap]} .
	\end{tikzcd}\]
\end{definition}

容易验证它们都是良定义的. 结合 (7.7.3) 和引理 7.7.9 可见对于所有右 $A$-模 $N$ 和右 $B$-模 $M$, 伴随对 (7.7.6) 的单位态射 $\eta_N$ 和余单位态射 $\epsilon_M$ 满足
\[\begin{tikzcd}[row sep=tiny]
	N \arrow[r, "\eta_N"] & \Hom_{\cated{Mod}B}(P, N \dotimes{A} P) & N \dotimes{A} P \dotimes{B} P^\vee \arrow[l, "\sim"'] \\
	x \arrow[mapsto, r] & {[p \mapsto x \otimes p]} & x \otimes \mathrm{coev}(1) \arrow[mapsto, l] \\
	M \dotimes{B} P^\vee \dotimes{A} P \arrow[r, "\sim"] & \Hom_{\cated{Mod}B}(P, M) \dotimes{A} P \arrow[r, "\epsilon_M"] & M \\
	x \otimes \lambda \otimes p \arrow[mapsto, r] & x \lambda(\cdot) \otimes p \arrow[mapsto, r] & x\lambda(p) \arrow[equal, d] \\
	& & (\identity_M \otimes \mathrm{ev})(x \otimes \lambda \otimes p).
\end{tikzcd}\]
因此我们无妨作等同
\begin{align*}
	\eta_N = \identity_N \otimes \mathrm{coev}: \; & N \to N \dotimes{A} P \dotimes{B} P^\vee , \\
	\epsilon_M = \identity_M \otimes \mathrm{ev}: \; & M \dotimes{B} P^\vee \dotimes{A} P \to M.
\end{align*}
\par\medskip\noindent\textit{作者原文：chapter7.tex，行 1968--1975.}\par
\def\CorpusNumber{7.7.15}
\begin{lemma}\label{prop:fg-module-cat}
	设 $R$ 为环, $N$ 为右 $R$-模, 则 $N$ 是有限生成的当且仅当以下性质成立: 对 $N$ 的任一族子对象 $(N_i)_{i \in I}$, 若 $\sum_{i \in I} N_i = N$, 则存在有限子集 $I_0 \subset I$ 使得 $\sum_{i \in I_0} N_i = N$.
\end{lemma}
\begin{proof}
	对于``仅当''方向, 设 $x_1, \ldots, x_n$ 为 $N$ 的一族生成元, 每个 $x_j$ 皆包含于某个 $\sum_{i \in I_j} N_i$, 其中 $I_j \subset I$ 有限; 取 $I_0 = \bigcup_{j=1}^n I_j$ 便是. 对于``当''的方向, 考虑 $N = \sum_{x \in N} xR$.
\end{proof}

现在回到定义--命题 7.7.4 介绍的森田等价. 我们只论右模情形.
\subsection*{标准先修与保留说明（整理者）}
通常的伴随、锥与（余）极限、Abel 范畴、投射对象和生成元定义为标准先修。原书编号引用在本文件中保持原编号。森田分类中 Eilenberg--Watts 定理 7.7.2 为具名标准先修；本文件保留其陈述，以及实际用到的有限投射对偶、评价与余评价构造。原书命题 2.8.15 为投射生成元的忠实正合刻画；A.2.2 为紧对象定义（已在局部背景中复述）。原文的简短例行验证及任何已存在的笔误均原样保留，未作数学修改。
\begin{thebibliography}{Li1}
\bibitem{Li1} 李文威，《代数学方法》第一卷：基础架构，高等教育出版社。原书引用保留原章节和结果编号；\url{https://www.wwli.asia/mathematics/AlJabr/}.
\end{thebibliography}
\end{document}
